%!TeX program = xelatex
\documentclass{SYSUReport}
\usepackage{array,tabularx,placeins,caption,float}
\usepackage{amsmath,amssymb}
\usepackage{tikz}
\usepackage{pgfplots}
\pgfplotsset{compat=1.16}
\usepackage{listings}
\IfFontExistsTF{Consolas}{\setmonofont{Consolas}}{}
\lstdefinestyle{pycode}{
  language=Python,
  basicstyle=\ttfamily\footnotesize,
  keywordstyle=\color{myteal}\bfseries,
  commentstyle=\color{gray!80!black},
  stringstyle=\color{orange!60!black},
  numbers=left,
  numberstyle=\tiny\color{gray},
  numbersep=6pt,
  frame=single,
  rulecolor=\color{gray!45},
  backgroundcolor=\color{gray!6},
  breaklines=true,
  breakatwhitespace=true,
  columns=fullflexible,
  keepspaces=true,
  showstringspaces=false,
  tabsize=4,
  xleftmargin=1.6em,
  xrightmargin=0.2em,
  aboveskip=6pt,
  belowskip=8pt,
}
\lstset{style=pycode}
\captionsetup{font=small,labelfont=bf,skip=7pt}
\setlength{\headheight}{16pt}
\setlength{\parskip}{0.2em}
\setlength{\tabcolsep}{6pt}
\renewcommand{\arraystretch}{1.22}
\linespread{1.18}
\IfFontExistsTF{SimHei}{\setCJKfamilyfont{zhhei}[AutoFakeBold=2]{SimHei}}{}
\clubpenalty=10000
\widowpenalty=10000
\displaywidowpenalty=10000
\raggedbottom
\IfFontExistsTF{Times New Roman}{\setmainfont{Times New Roman}}{}
\headl{数字图像处理}
\headc{}
\headr{基于颜色的人脸分割}
\lessonTitle{数字图像处理实验报告}
\reportTitle{基于颜色的人脸分割}
\stuname{}
\stuid{}
\inst{}
\major{}
\date{2026 年 11 月 2 日}
\hypersetup{pdftitle={基于颜色的人脸分割},pdfauthor={},pdfsubject={颜色空间、肤色统计、区域后处理与形态学补全}}
\begin{document}
\hypersetup{pageanchor=false}
\cover
\thispagestyle{empty}
\clearpage
\pagenumbering{Roman}
\hypersetup{pageanchor=true}

\begin{abstract}
本报告研究颜色空间中的人脸区域分割。实验使用 162 张带像素级标签的图像，按 98/32/32 固定划分；训练集只取皮肤与鼻子像素估计 Cr/Cb 的均值与标准差，验证集比较阈值系数与后处理，测试集只用于固定方案的评价。独立测试中，$k=1.25$ 的 Cr/Cb 阈值配合最大连通域与填孔取得 F1 70.69\%、IoU 54.67\%，比纯阈值基线提高 8.03 和 9.04 个百分点，收益主要来自减少误检。

针对独立测试暴露的内部缺口，我们尝试在填孔前加入 $41\times41$ 膨胀、填孔后再腐蚀的补全流程。在同一批测试图上，召回提高 4.00 个百分点，精确率下降 2.80 个百分点，F1 基本不变（70.70\%）。新增前景中只有 35.45\% 属于脸部，误检主要来自头发。这一结果表明，形态学只能移动精确率—召回率工作点，不能改变颜色模型的语义上限：耳朵、脖子与脸部同色，纯颜色判别无法区分；头发与背景的混淆也不能靠补洞解决。

报告据此讨论方法的有效条件、失效层次与适用范围，并把两次实现放在同一张精确率—召回率平面上比较。主要结论是：在受控条件下，颜色模型可以给出可用的粗区域；它的瓶颈不在阈值或核大小，而在颜色本身不包含“这是脸”的语义。

\noindent\textbf{关键词：}人脸分割；肤色检测；YCbCr；形态学补全；精确率—召回率
\end{abstract}

\clearpage
\tableofcontents
\clearpage
\pagenumbering{arabic}

\section{问题与路线}\label{sec:intro}

\subsection{任务与边界}

课程要求在一个颜色空间内对亚洲人脸进行分割，并对多幅人脸图像验证。我们把它落实为像素级的人脸区域分割：输入一张 $512\times512$ 的彩色图像，输出与原图同尺寸的二值掩膜，1 表示脸部，0 表示其他。评价的正类包括皮肤、鼻子、双眼、双眉、口腔和上下唇；眼镜作为忽略区域，不参加任何混淆计数。

有两个边界需要先说明。第一，肤色掩膜不等于人脸检测，更不等于身份识别。掩膜可能包含脖子、耳朵等皮肤区域，也不保证画面中只有一张脸；本报告不评价识别性能。第二，数据中的“亚洲人脸”来自实验者的外观筛选，没有官方人群属性标签；报告中的结论只对应这批图像及其固定划分，不外推到“所有亚洲人脸”。

\subsection{三个问题}

全文围绕三个问题展开。第一，这批人脸的肤色在 Cr/Cb 上是否足够集中，均值与标准差能否概括？第二，只靠颜色加两个空间规则，F1 能做到多少，提升来自哪一步？第三，它在什么图像条件下还成立，失败是参数问题还是模型问题？

对应地，我们经历了两个实验阶段的递进探索。实验三围绕纯颜色与基础连通域规则做了三次判断：第一次是颜色范围（$k$ 取多少），第二次是空间结构（最大连通域与填孔值多少），第三次是无约束补全强度（在填孔前加入 $41\times41$ 膨胀、填孔后再腐蚀能否补回内部缺口）。实验四则针对实验三暴露的“过早取最大连通域误删脸部碎片”与“头发脖颈粘连无法靠无约束形态学分离”两类核心瓶颈，引入基于灰度黑帽变换与双窗先验的双眼位置识别，由双眼几何关系生成正中与侧向内外嵌套双椭圆，结合低支持度样本放宽肤色阈值（$k=2.00$）重选及分区形态学连接，突破了纯颜色模型的等值线限制。

\subsection{证据状态与章节安排}

初版方案在固定配置后只运行一次测试，得到本报告唯一的独立测试结果。实验三的 $41\times41$ 补全尝试来自测试图 00570 的误差，因此在同一批测试图上的数字按开发后再评价报告；实验四的双眼定位与双椭圆约束形态学优化则严格固定在 32 张验证集（开发集）上进行多轮消融与验证，各阶段结果在表与图中均明确标注所属集合与证据状态，不相互混用。

第~\ref{sec:data}~章说明数据、标签与评价口径；第~\ref{sec:model}~章给出颜色模型与实验三初版方案；第~\ref{sec:test}~章报告独立测试并归纳三类失败；第~\ref{sec:completion}~章接第三类失败，做无约束形态学补全尝试并再次评价；第~\ref{sec:boundary}~章讨论纯颜色模型的能力边界与适用范围；第~\ref{sec:lab4}~章系统报告实验四的双眼定位、双椭圆空间约束、侧向椭圆候选、低支持放宽阈值重建与分区形态学连接优化；第~\ref{sec:conclusion}~章总结全文。

在独立测试集上，实验三初版 F1 为 70.69\%，$41\times41$ 补全后为 70.70\%，说明无空间锚点的形态学只能改变召回与精确率的分配；而在验证集上，实验四将全局 F1 从实验三的 72.70\% 大幅提升至 82.89\%（IoU 从 57.11\% 提升至 70.79\%），证明引入五官位置与双椭圆几何先验能够实质性跨越纯颜色分割的性能瓶颈。

\FloatBarrier
\section{数据与评价口径}\label{sec:data}

\subsection{数据来源与筛选}

实验使用 Face\_segment 数据集及其原始多类别像素级标签。我们从 \texttt{test/real\_image/2} 目录中筛选出外观符合亚洲人脸特征的样本，并保留对应的像素标签。初选共得 163 对图像与标签；在配对、尺寸与标签取值核查中，发现样本 00583 与 00584 的原图完全一致而标签存在差异，经核对保留 00583、剔除 00584，最终确立 162 对有效样本。所有原图均为 $512\times512$ 的 RGB 彩色图像，标签为单通道类别图，包含 0--18 共 19 个语义类别。通过对原图、彩色编码标签与半透明叠加图的逐张核查，确认图像与标注配对准确、边界清晰。

\subsection{数据集划分与去重检查}

我们将 162 对样本按编号排序后，使用固定随机种子 42 按约 $6:2:2$ 的比例划分为训练集 98 张、验证集 32 张与测试集 32 张。在划分前，采用 64 位感知哈希（dHash，汉明距离阈值 $\le 8$）对全集进行近重复筛查，并对筛出的 35 对候选图像逐一复核，确认不存在同一照片裁剪或缩放重复。训练集用于统计肤色模型参数与五官几何先验，验证集用于算法消融与超参数选择，测试集用于检验基础流程的泛化表现。

\subsection{两套标签映射：建模用皮肤，评价用整脸}

本实验在“肤色统计建模”与“完整脸部区域评价”两个环节采用不同的标签子集：肤色建模仅提取皮肤和鼻子像素；而在评价完整脸部分割效果时，正类除皮肤和鼻子外，还包含双眼、双眉、口腔和上下唇。眼镜区域设为忽略类别，其余类别均归为负类。完整映射关系见表~\ref{tab:labels}。

\begin{table}[H]
\centering\small
\caption{原始标签在颜色统计与脸部评价中的映射规则}\label{tab:labels}
\begin{tabularx}{\textwidth}{l l X}
\toprule
用途 & 标签编号 & 对应区域 \\
\midrule
肤色统计 & 1、2 & 皮肤、鼻子 \\
脸部评价正类 & 1、2、4--7、10--12 & 皮肤、鼻子、双眼、双眉、口腔、上下唇 \\
评价忽略区域 & 3 & 眼镜 \\
脸部评价负类 & 其余类别 & 背景、双耳、头发、帽子、饰品、脖颈、衣服 \\
\bottomrule
\end{tabularx}
\end{table}

这一差异化设计符合人脸图像的物理特性：眼睛、眉毛和嘴部的颜色通常偏离皮肤色度分布，若将其纳入肤色统计会污染均值与方差估计；但完整的人脸掩膜应当涵盖五官区域，因此颜色阈值筛选在五官位置留下的天然空洞，恰好成为后续连通域填孔与形态学修补的作用对象。实现上通过集合包含关系生成二值真值掩膜与有效评测掩膜，见代码~\textbf{2-1}。

\noindent\textbf{代码 2-1}\quad 脸部真值与有效区域的生成（dataset.py）
\begin{lstlisting}
POSITIVE_IDS = [1, 2, 4, 5, 6, 7, 10, 11, 12]
IGNORE_IDS = [3]

def binary_labels(labels):
    face = np.isin(labels, POSITIVE_IDS).astype(np.uint8)
    valid = np.logical_not(np.isin(labels, IGNORE_IDS)).astype(np.uint8)
    return face, valid
\end{lstlisting}

代码中 \texttt{face} 标记完整脸部正类，\texttt{valid} 排除眼镜区域；评价时仅在 \texttt{valid} 为 1 的像素上统计混淆矩阵。

\begin{figure}[H]
\centering
\includegraphics[width=\textwidth]{part_a_label_review_00546.png}
\caption{样本 00546 的原图、原始多类别标签及半透明叠加图。}
\label{fig:label}
\end{figure}

图~\ref{fig:label} 展示了原始标签对皮肤、鼻子、眼睛、眉毛和嘴部的精细划分，同时可以看到头发和脖颈虽与脸部紧邻，但在评测定义中属于负类。明确这一语义边界有助于区分“面部内部五官漏检”与“面部外围头发脖颈误检”。

\subsection{评价指标体系}

在测试集 32 张图像共 8,388,608 个有效像素中，脸部正类像素为 2,150,722 个，占 25.64\%，背景与非脸部负类占 74.36\%。由于正负类比例显著失衡，若将全图所有像素均判为背景，像素准确率（Accuracy）仍高达 74.36\%，但此时 F1 为 0；反之，若将全图均判为脸部，F1 仅为 40.80\%（见表~\ref{tab:trivial}）。

\begin{table}[H]
\centering\small
\caption{平凡预测与基础方法在测试集上的指标对照，单位为百分比}\label{tab:trivial}
\begin{tabular}{lrrr}
\toprule
预测方式 & Accuracy & Precision & F1 \\
\midrule
全判为背景 & 74.36 & -- & 0 \\
全判为脸部 & 25.64 & 25.64 & 40.80 \\
纯阈值基线（$k=1.50$） & 76.53 & 52.91 & 62.66 \\
实验三基础方案 & 83.94 & 66.43 & 70.69 \\
\bottomrule
\end{tabular}
\end{table}

因此，本报告以 F1 与交并比（IoU）作为核心综合指标，辅以精确率（Precision）和召回率（Recall）刻画误检与漏检之间的权衡，Accuracy 仅作参考。多图汇总采用两种口径：一是先累加全部图像的 TP、FP、TN、FN 再计算整体指标（全局指标，Global），二是对每张图像单独计算指标后再取算术平均（逐图均值，Per-image Mean）。正文默认报告全局指标，并在消融与稳定性分析中同时列出逐图均值。各指标定义如式~\eqref{eq:metrics} 所示：
\begin{equation}
\begin{aligned}
\mathrm{Precision}&=\frac{TP}{TP+FP}, &
\mathrm{Recall}&=\frac{TP}{TP+FN},\\
\mathrm{F1}&=\frac{2TP}{2TP+FP+FN}, &
\mathrm{IoU}&=\frac{TP}{TP+FP+FN}.
\end{aligned}
\label{eq:metrics}
\end{equation}

\FloatBarrier
\section{颜色模型与初版方案}\label{sec:model}

\subsection{颜色空间与丢弃 Y 的代价}

YCbCr 把亮度与色度分开。按 ITU-R BT.601，RGB 到 YCbCr 的常用形式为
\begin{equation}
\begin{aligned}
Y &= 0.299R + 0.587G + 0.114B,\\
Cr &= 0.713(R-Y) + 128,\\
Cb &= 0.564(B-Y) + 128.
\end{aligned}
\label{eq:ycbcr}
\end{equation}
OpenCV 的整数实现与上式略有差异，但通道顺序与含义一致：Y、Cr、Cb，8 位编码 0--255。实现按 BGR 读图，用 \texttt{COLOR\_BGR2YCrCb} 转换，不做缩放、白平衡或颜色增强。

只用 Cr、Cb 的理由是肤色在色度平面上比在 RGB 中更集中，亮度变化主要落在 Y 上\cite{hsu2002face,vezhnevets2003survey}。代价是丢掉了 Y 中的明暗线索：深色头发和皮肤在色度上可能接近，但在亮度上差别明显。00718 后来正暴露出这一代价；这是当时就知道、但接受了的取舍。

\subsection{训练肤色统计}

从 98 张训练图取皮肤与鼻子像素，共 $5{,}838{,}561+548{,}991=6{,}387{,}552$ 个。所有像素等权，不随机抽样，也不先求各图均值再平均；脸部面积大的图像贡献更多像素。对通道 $c\in\{Cr,Cb\}$，均值与总体方差为
\begin{equation}
\mu_c=\frac{1}{N}\sum_{i=1}^{N}c_i,\qquad
\sigma_c^2=\frac{1}{N}\sum_{i=1}^{N}(c_i-\mu_c)^2,
\label{eq:stats}
\end{equation}
其中 $N=6{,}387{,}552$。计算结果如表~\ref{tab:skin} 上半部分。两通道标准差都约为 9.4，离散尺度接近；这并不表示分类能力相同，因为统计里没有负类。

\begin{table}[H]
\centering\small
\caption{训练肤色统计与阈值覆盖}\label{tab:skin}
\begin{tabular}{lrrr}
\toprule
通道 & 均值 & 总体方差 & 标准差 \\
\midrule
Cr & 151.1495 & 87.8810 & 9.3745 \\
Cb & 109.1102 & 89.9380 & 9.4836 \\
\bottomrule
\end{tabular}

\vspace{6pt}
\begin{tabular}{lrrr}
\toprule
阈值系数 $k$ & 1.00 & 2.00 & 3.00 \\
范围内训练肤色像素占比 & 59.17\% & 91.86\% & 99.36\% \\
\bottomrule
\end{tabular}
\end{table}

\noindent\textbf{代码 3-1}\quad 训练肤色的直方图统计（modeling.py）
\begin{lstlisting}
mask = (labels == 1) | (labels == 2)
ycrcb = cv2.cvtColor(bgr, cv2.COLOR_BGR2YCrCb)
cr = ycrcb[:, :, 1][mask].astype(np.int64)

counts = np.bincount(cr, minlength=256)
n = counts.sum()
values = np.arange(256, dtype=np.float64)
mean = np.dot(values, counts) / n
std = np.sqrt(np.dot((values - mean) ** 2, counts) / n)
\end{lstlisting}

这段代码先在每张训练图上取皮肤与鼻子像素，累计 256 维直方图，再跨图求和。因为统计对象是直方图而不是逐图像均值，每个像素的权重相同；脸部面积大的图像自然贡献更多像素。方差除以 $N$ 而不是 $N-1$，对应总体方差；$\sigma$ 由方差开方得到，阈值中使用的是 $\sigma$。

联合分布见图~\ref{fig:joint}。高频区域呈倾斜椭圆，按像素频数计算的相关系数约为 $-0.7644$，协方差为 $-67.95$：Cr 较大的像素往往对应较小的 Cb。若分别限制两个通道，得到的是轴对齐矩形；它无法贴合倾斜的分布，四角会包含肤色之外的组合，倾斜方向上的肤色又可能被切掉。

\begin{figure}[H]
\centering
\includegraphics[width=0.74\textwidth]{part_a_joint_distribution.png}
\caption{训练肤色像素的 Cr/Cb 联合分布。横轴为 Cb，纵轴为 Cr；颜色表示对应整数取值组合的像素数，采用对数刻度。}
\label{fig:joint}
\end{figure}

这个数还有一个用途：估计阈值覆盖了多少训练肤色。表~\ref{tab:skin} 下半部分给出 $k=1,2,3$ 的覆盖比例。$k=1$ 时矩形只覆盖 59.17\%，$k=2$ 才到 91.86\%。也就是说，最终选的 $k=1.25$ 不是“肤色范围”，而是一个偏精确率的切片：$k=1$ 时就已经只覆盖 59.17\% 的训练肤色，$k=1.25$ 也远没有覆盖全部。它同时解释了两件事：初版的召回为什么上不去，以及低亮度图像为什么一收紧就掉得厉害。

\subsection{判决规则与验证选参}

对每个像素，同时满足式~\eqref{eq:box} 时记为候选前景：
\begin{equation}
M(x,y)=
\begin{cases}
1, & |Cr(x,y)-\mu_{Cr}|\leq k\sigma_{Cr}\ \text{且}\ |Cb(x,y)-\mu_{Cb}|\leq k\sigma_{Cb},\\
0, & \text{其他情况}.
\end{cases}
\label{eq:box}
\end{equation}
实现中先转浮点再作差，包含边界；不对每张图像重新估计参数。

颜色候选还需要空间处理。脸部真值包括眼睛、眉毛、口腔与嘴唇，这些位置未必满足肤色范围，容易形成空洞；与脸部分离的类肤色背景则可以通过连通域筛选排除。最终方案先按 8 连通保留面积最大的前景块，再填充不与图像外部背景 4 连通的内部空洞。

\noindent\textbf{代码 3-2}\quad 阈值判决与填孔（detection.py）
\begin{lstlisting}
converted = cv2.cvtColor(image_bgr, cv2.COLOR_BGR2YCrCb).astype(np.float64)
cr, cb = converted[:, :, 1], converted[:, :, 2]
mask = ((np.abs(cr - mean_cr) <= k * std_cr) &
        (np.abs(cb - mean_cb) <= k * std_cb)).astype(np.uint8)

def fill_holes(mask):
    padded = np.pad(mask.astype(np.uint8), 1)
    cv2.floodFill(padded, None, (0, 0), 2, flags=4)
    return (padded[1:-1, 1:-1] != 2).astype(np.uint8)
\end{lstlisting}

这段代码有两个细节。第一，相减前先转成浮点，避免 \texttt{uint8} 回绕；比较使用 \texttt{<=}，边界像素算作候选。第二，填孔先在外围补一圈零，再从 \texttt{(0,0)} 出发按 4 连通把外部背景标成 2；没有被标到的零像素就是被前景围住的洞，最后统一置 1。补零保证即使目标贴着图像边缘，外部种子仍然有效。

验证集 32 张，比较 9 个 $k$（1 到 3，步长 0.25）与 4 种处理（无、填孔、最大区域加填孔、闭运算后最大区域加填孔），共 36 组。主指标是全局 F1；平局先选更简单的处理，再选较小的 $k$。计划在运行前写入配置文件。

\begin{figure}[H]
\centering
\includegraphics[width=0.96\textwidth]{validation_curve.png}
\caption{四种处理方式在验证集上的参数曲线。横轴为颜色阈值系数 $k$，纵轴为全局 F1。}
\label{fig:curve}
\end{figure}

图~\ref{fig:curve} 有两个特征。第一，1.25 到 1.5 之间几乎是平台，两者相差不到 0.5 个百分点；最终选 1.25，一部分原因是平局规则，一部分原因是更小的 $k$ 在测试前更保守。第二，$k$ 超过 1.5 后 F1 断崖式下降：颜色范围一放宽，头发和背景就被并进最大连通域，F1 从 72 降到 65、再到 59。失效是结构性的，不是渐进的精度损失。

在相同 $k=1.25$ 下比较四种后处理，结果如表~\ref{tab:post}。填孔把 F1 从 56.13 拉到 69.10，是收益最大的一步；最大连通域再加 3.60 个点；闭运算反而低 0.71 个点。我们原本更看好闭运算，觉得它才像“把眉眼补上”，但它在验证集上不如填孔。原因后来在补全实验里才想清楚：闭运算先改变连通性，会把脸和脖子、头发粘在一起，最大区域更容易选错；填孔只填被围住的洞，不动外轮廓。

\begin{table}[H]
\centering\small
\caption{相同颜色阈值下四种后处理的验证结果，单位为百分比}\label{tab:post}
\begin{tabular}{lrrrr}
\toprule
后处理 & Precision & Recall & F1 & IoU \\
\midrule
无处理 & 52.43 & 60.39 & 56.13 & 39.01 \\
填孔 & 58.34 & 84.71 & 69.10 & 52.78 \\
最大连通域加填孔 & 64.86 & 82.70 & 72.70 & 57.11 \\
闭运算后最大区域加填孔 & 63.39 & 83.30 & 71.99 & 56.24 \\
\bottomrule
\end{tabular}
\end{table}

最终固定 $k=1.25$、最大连通域加填孔，记为初版方案。它在验证集上的 P/R/F1/IoU 为 64.86/82.70/72.70/57.11。这个配置随后只运行一次测试，结果见第~\ref{sec:test}~章。

\subsection{设计取舍}

表~\ref{tab:tradeoff} 列出当时主动放弃的几种做法。它们不是事后补的理由，而是当时的设计边界；第~\ref{sec:test}~章与第~\ref{sec:completion}~章的失败会逐一回到这些边界上。

\begin{table}[H]
\centering\small
\setlength{\tabcolsep}{4pt}
\caption{设计取舍与后来看到的代价}\label{tab:tradeoff}
\begin{tabular}{p{3.1cm}p{5.0cm}p{6.2cm}}
\toprule
放弃的做法 & 原因 & 后来看到的代价 \\
\midrule
使用 Y 通道 & 让阈值只管颜色，减少曝光影响 & 深色头发与皮肤难分（00718） \\
逐图自适应阈值 & 保持一个全局模型，避免每图调参 & 偏色图像无法自救（00609） \\
椭圆或马氏距离 & 实现简单、结果可解释 & 矩形四角包含非肤色；$-0.76$ 的相关性未利用 \\
深度分割模型 & 课程目标与可解释性 & 语义鸿沟只能靠后处理绕开 \\
多人处理 & 数据中没有多人标注 & 最大连通域只保留一个主体 \\
\bottomrule
\end{tabular}
\end{table}

其中“椭圆或马氏距离”值得多说一句。RGB 规则更简单，但三个通道都混有亮度信息\cite{kovac2003skin}；椭圆边界不需要更换颜色空间，只需在均值、标准差之外增加一个相关系数，就能利用 $-0.76$ 的联合关系。本实验没有实现它，因此不能给出它可能带来的提升；这里只把它作为由数据直接支持的一个判断记录下来。

\FloatBarrier
\section{初版方案的独立测试}\label{sec:test}

\subsection{测试设置}

测试使用固定划分中的 32 张图，每张 $512\times512$，共 8,388,608 个有效像素。检测只读原图、训练统计和固定配置；评价程序再读取保存的预测，与独立生成的真值对齐。测试集没有参与选参，也没有在测试后修改配置。当前测试图中没有眼镜像素，忽略逻辑保留但未被触发。

评价只做像素计数，不调用检测函数内部的状态；核心循环见代码~\textbf{4-1}。

\noindent\textbf{代码 4-1}\quad 有效区域内的混淆计数与指标（evaluation.py）
\begin{lstlisting}
valid = (np.ones(gt.shape, dtype=bool) if valid_mask is None
         else valid_mask.astype(bool))
not_pred, not_gt = np.logical_not(pred), np.logical_not(gt)
tp = np.count_nonzero(pred & gt & valid)
fp = np.count_nonzero(pred & not_gt & valid)
tn = np.count_nonzero(not_pred & not_gt & valid)
fn = np.count_nonzero(not_pred & gt & valid)

precision = tp / (tp + fp) if tp + fp else None
recall    = tp / (tp + fn) if tp + fn else None
f1        = 2 * tp / (2 * tp + fp + fn) if (2 * tp + fp + fn) else None
iou       = tp / (tp + fp + fn) if (tp + fp + fn) else None
\end{lstlisting}

这段代码用逻辑与同时完成“预测为正”和“真值为正”的筛选：\texttt{pred} 是预测掩膜，\texttt{gt} 是脸部真值，二者都是 0/1 数组；\texttt{valid} 把眼镜像素排除在四个计数之外。分母为零时返回 \texttt{None}，不做静默填零。汇总时先把所有图像的 TP/FP/TN/FN 相加再计算 global，逐图平均则单独对每张图的指标做等权平均；两种口径在报告中始终分开标注。

\subsection{总体结果}

总体指标与混淆计数分别见表~\ref{tab:test} 和表~\ref{tab:testcounts}。

\begin{table}[H]
\centering\small
\caption{测试集总体指标，单位为百分比}\label{tab:test}
\begin{tabular}{lrrrr}
\toprule
方案 & Precision & Recall & F1 & IoU \\
\midrule
纯阈值基线 & 52.91 & 76.82 & 62.66 & 45.63 \\
初版方案 & 66.43 & 75.55 & 70.69 & 54.67 \\
变化，百分点 & +13.52 & $-1.27$ & +8.03 & +9.04 \\
\bottomrule
\end{tabular}
\end{table}

\begin{table}[H]
\centering\small
\caption{测试集混淆计数}\label{tab:testcounts}
\begin{tabular}{lrrrr}
\toprule
方案 & TP & FP & TN & FN \\
\midrule
纯阈值基线 & 1,652,227 & 1,470,496 & 4,767,390 & 498,495 \\
初版方案 & 1,624,818 & 821,271 & 5,416,615 & 525,904 \\
\bottomrule
\end{tabular}
\end{table}

初版相对基线：F1 提高 8.03 个百分点，IoU 提高 9.04 个百分点；Precision 提高 13.52，Recall 下降 1.27。计数上，FP 减少 649,225，FN 增加 27,409。收益主要来自减少误检，而不是找回漏检。

初版相对基线的变化主要是“删掉不该选的”，而不是“把漏掉的找回来”。召回基本不动，精确率大幅上升；这也解释了为什么初版的 F1 停在 70\% 附近——它把颜色候选修剪得更干净，但没有增加新的判别信息。

\subsection{提升来自哪一步}

把流程拆成四个阶段，可以看到每一步拿什么换什么，结果如表~\ref{tab:stage}。

\begin{table}[H]
\centering\small
\setlength{\tabcolsep}{4pt}
\caption{固定流程各阶段的测试结果，指标单位为百分比}\label{tab:stage}
\begin{tabular}{lrrrrr}
\toprule
阶段 & Precision & Recall & F1 & FP & FN \\
\midrule
基线阈值，$k=1.50$ & 52.91 & 76.82 & 62.66 & 1,470,496 & 498,495 \\
颜色阈值，$k=1.25$ & 56.59 & 65.24 & 60.61 & 1,076,315 & 747,523 \\
保留最大连通域 & 64.73 & 63.36 & 64.04 & 742,515 & 787,938 \\
填孔后的结果 & 66.43 & 75.55 & 70.69 & 821,271 & 525,904 \\
\bottomrule
\end{tabular}
\end{table}

三步的交换率不同。$k$ 从 1.50 收到 1.25，去掉 394,181 个误检像素，也丢掉 249,028 个正确像素，约合 1 个 TP 换 1.6 个 FP；最大连通域去掉 333,800 个误检，只损失 40,415 个正确像素，约合 1 个 TP 换 8.3 个 FP；填孔补回 262,034 个正确像素，同时增加 78,756 个误检，约合 1 个 FP 换 3.3 个 TP。

单看 $k$ 收紧这一步，F1 从 62.66 降到 60.61，是亏的；它是在给后面的最大连通域让路。最大连通域是最划算的一步；填孔再把内部补回来。三步不能分开评价：阈值、区域和填孔之间存在配合关系，单独调其中一步会得到不同的“最优”。

\subsection{误检从哪里来}

表~\ref{tab:errors} 把最终误检按原始类别归类，并给出填孔前后两组正类的组内召回。

\begin{table}[H]
\centering\small
\caption{最终误检的语义来源与五官召回}\label{tab:errors}
\begin{tabular}{lrrr}
\toprule
负类 & 基线 FP & 最终 FP & 占最终 FP \\
\midrule
头发 & 772,905 & 497,516 & 60.58\% \\
脖颈 & 248,769 & 192,139 & 23.40\% \\
背景 & 369,581 & 93,639 & 11.40\% \\
衣服 & 58,056 & 20,420 & 2.49\% \\
双耳 & 17,665 & 14,184 & 1.73\% \\
其他 & 3,520 & 3,373 & 0.41\% \\
\bottomrule
\end{tabular}

\vspace{6pt}
\begin{tabular}{lrr}
\toprule
正类分组 & 最大区域后 & 填孔后 \\
\midrule
皮肤与鼻子 & 66.43 & 76.27 \\
眼、眉、口腔与嘴唇 & 32.61 & 68.32 \\
\bottomrule
\end{tabular}
\end{table}

头发和脖颈合计占最终误检的 83.98\%。头发容易形成类肤色候选；脖颈则与脸部颜色相近，又常常连在同一前景块上，最大连通域无法把它切开。所以剩余问题不只来自分离的背景，也来自与脸部同色或相连的区域。

其中脖子和耳朵本身是皮肤，这是方法的语义上限。测试集中脖子与双耳合计约 37.4 万像素，而脸部正类约 215 万像素。即使颜色判决完全正确、不再产生任何头发或背景误检，只要结构仍是“最大连通域加填孔”，这部分像素仍会被计入前景，精确率的上限约为 85\%。当前精确率 66.43\%，离上限还有约 19 个百分点；即使把头发和背景误检全部消除，纯颜色模型也无法达到更高的精确率。这个计算假设所有肤色像素都被判为正类，实际方法可能因连通性排除部分耳朵，因此 85\% 是量级上的上限，不是严格界。

填孔确实恢复了五官：眼、眉、口腔与嘴唇的组内召回从 32.61\% 提高到 68.32\%，补回 69,662 个正确像素；皮肤与鼻子组从 66.43\% 提高到 76.27\%。方法并不知道那些像素是眼睛或嘴，它只是利用了它们被脸部轮廓包围这一拓扑事实。

\subsection{逐图稳定性与亮度条件}

逐图指标与亮度分组见表~\ref{tab:perimage}。

\begin{table}[H]
\centering\small
\caption{逐图指标、边界质量与亮度分组}\label{tab:perimage}
\begin{tabular}{lrr}
\toprule
统计项目 & 基线 & 初版 \\
\midrule
逐图 F1 平均 & 0.628 & 0.681 \\
逐图 F1 中位数 & 0.664 & 0.778 \\
逐图 F1 最小值 & 0.273 & 0.031 \\
逐图 F1 最大值 & 0.907 & 0.930 \\
Boundary IoU 平均 & 0.153 & 0.233 \\
\bottomrule
\end{tabular}

\vspace{6pt}
\begin{tabular}{lrrr}
\toprule
亮度分组 & 图像数 & 基线 F1 & 初版 F1 \\
\midrule
较低 & 11 & 56.56 & 52.82 \\
中间 & 11 & 61.90 & 74.88 \\
较高 & 10 & 69.43 & 83.00 \\
\bottomrule
\end{tabular}
\end{table}

32 张图中有 27 张 F1 提高、5 张下降。中位数从 0.664 升到 0.778，但最小值从 0.273 降到 0.031：多数图像受益，少数图像变差，而且变差得更厉害。逐图 F1 几乎呈两段分布：11 张不低于 0.85，5 张低于 0.2，0.2 到 0.5 之间没有样本。这个形态与“方法依赖一个占主导的肤色连通区域”的解释一致；样本只有 32 张，形态也可能带有偶然性。

按图像配对重采样 10,000 次，总体 F1 差值的 95\% 区间为 $+3.57$ 至 $+12.00$ 个百分点；逐图平均 F1 差值的区间为 $-0.68$ 至 $+10.60$ 个百分点。后者跨过 0，说明改善集中在部分图像上；只报一个总体 F1 会掩盖这一点。

按真值脸部平均 $Y$ 排序，把测试图分成 11、11、10 张的较低、中间、较高亮度三组。较低亮度组的 F1 反而从 56.56 降到 52.82，正确像素减少约 9.6 万、漏检增加约 9.6 万；中、高亮度组则分别提高 13 和 13.6 个百分点。低亮度下肤色色度偏离训练分布，$k=1.25$ 把真肤色也切掉了，减少的误检不足以补偿。需要说明，亮度分组是事后按真值区域平均 $Y$ 划分的，只有 11 张图，且亮度与场景、肤色、妆容混杂；这里只能说明低亮度条件下风险更高，不能把下降单独归因于亮度。

另外采用 Boundary IoU 观察轮廓贴合程度\cite{cheng2021boundary}：使用预测和真值各自的内侧边界带，带宽取图像对角线的 2\%（当前为 14 像素），逐图平均后从 0.153 升到 0.233。区域重合改善的同时，边界也有所改善，但 0.233 仍然不高；掩膜看起来完整，边缘并不可信。

\begin{figure}[H]
\centering
\includegraphics[width=\textwidth]{sample_stability.png}
\caption{逐图 F1 比较与累积分布。左图每个点对应一张测试图，虚线上方表示初版更高；右图为 F1 累积分布。}
\label{fig:stability}
\end{figure}

图~\ref{fig:stability} 同时给出逐图变化和低分尾部。5 张退步图中，有 3 张的 Precision 是 100\%——00609、00583、00735 的召回分别只有 1.6\%、4.2\%、7.4\%。掩膜里留下的像素确实都在脸上，但整张脸几乎没被找到。只看“干不干净”，它们接近满分；按 F1 算，它们和全黑没有本质区别。

\subsection{三类失败}

把 32 张误差图逐张看过去，失败可以归成三类。

第一类是颜色候选不足。00609 在 $k=1.25$ 下只剩 5,685 个候选像素，分成 56 个连通块，最大块 1,063 像素；最终召回 1.6\%，F1 0.031。图像整体偏冷，肤色被推到训练范围之外；错误在颜色阶段就发生了，后面的区域选择和填孔只是没能救回来。

\begin{figure}[H]
\centering
\includegraphics[width=\textwidth]{test_00609.png}
\caption{00609 的原图、真值、初版预测与误差对照。颜色候选不足导致大面积漏检。}
\label{fig:00609}
\end{figure}

第二类是区域选择错误。00718 中，棕色头发形成 58,716 像素的最大连通块，其中正确像素只有 2,705 个；填孔又增加 110 个正确像素和 10,871 个误检像素，最终 F1 0.045。面积只能表示区域大小，不能表示区域属于什么。

\begin{figure}[H]
\centering
\includegraphics[width=\textwidth]{test_00718.png}
\caption{00718 的原图、真值、初版预测与误差对照。头发被选为主要候选，填孔扩大了误检。}
\label{fig:00718}
\end{figure}

第三类不同：主体大致正确，但眼睛、嘴部与外部背景连通，填孔补不回来。00570 是一个清楚的例子，嘴部有一条横向缺口；它的 F1 是 0.759，召回 0.631，主要缺失不是“选错”，而是“补不上”。前两类属于颜色与区域的边界，在现有框架内没有便宜的修法；第三类只涉及“如何把开口封住”，可以单独做一次实验。下一章从这里开始。

\FloatBarrier
\section{低支持度放宽肤色重选与分区形态学连接}\label{sec:low_support}

\subsection{内椭圆像素过少时的漏检现象}

在接通“双眼定位 + 内外双椭圆分区约束”的初版流程后，我们检查了 32 张验证图的结果。对于大多数在实验三中已经检出面部主体、只是外围粘连了头发或脖子的图像（例如 00611、00632），外椭圆裁剪与内椭圆闭运算已经能明显减少误检；但与此同时，有 4 张图像（00548、00562、00660、00733）的分割效果仍然很差，F1 停留在 $0.10\text{--}0.28$ 之间。

查看这 4 张图的中间掩膜可以发现，它们的共同问题是：在实验三 $k=1.25$ 阈值和最大连通域筛选之后，真正落入内椭圆 $E_{\text{in}}$ 的前景像素非常少（占内椭圆面积的比例 $\rho_{\text{in}}$ 都低于 $0.18$，例如 00562 在内椭圆里只剩下嘴角附近的一小块，00660 只选中了右下颌与脖子）。当内椭圆里几乎没有足够的初始肤色像素时，即使在内椭圆中做 $11\times11$ 闭运算和填孔，也没有东西可以连接，整张内脸依然是大片空白。

\subsection{内椭圆放宽肤色阈值重选与受限重建}

为什么这几张图在 $k=1.25$ 下落入内椭圆的像素这么少？回到第~\ref{sec:intro}~章提到的训练集肤色统计可以看到，$k=1.25$ 的矩形阈值本来只覆盖了约六到七成的训练肤色像素；当图像整体偏暗、受冷光影响或面部存在较大阴影时，真实脸部皮肤的 Cr、Cb 值很容易落到 $k=1.25$ 之外，但在 $k=2.00$（训练集肤色覆盖率达 91.86\%）下其实仍能被检出。

在实验三中，我们不敢直接把全图阈值放宽到 $k=2.00$，因为全图放宽会把大片头发和背景一起变成前景。但现在已经有了由双眼定出的内椭圆 $E_{\text{in}}$，情况就不同了：如果只在内椭圆像素明显不足的图像上，并且只允许在内椭圆 $E_{\text{in}}$ 内部使用放宽后的肤色阈值重选像素，外围背景和远端头发就不会被放进来。

按照这一思路，我们将“核心种子提取”和“低支持度放宽阈值重建”整理为以下三个步骤：
\begin{enumerate}
\item \textbf{提取高可信核心种子并计算内椭圆支持率：}设实验三输出的掩膜为 $M_{\text{Lab3}}$，先用外椭圆裁剪得到 $M_{\text{local}} = M_{\text{Lab3}} \cap E_{\text{out}}$，并在其中提取与内椭圆 $E_{\text{in}}$ 重叠像素数不少于 $N_{\text{core}}=24$ 的 8 连通分量，合记为 $M_{\text{sel}}$；若没有任何分量满足条件，则退一步保留 $M_{\text{local}}$ 中的最大连通块。将 $M_{\text{sel}}$ 落在内椭圆中的部分记为核心种子 $S = M_{\text{sel}} \cap E_{\text{in}}$，并计算内椭圆支持率
\begin{equation}
\rho_{\text{in}} = \frac{|S|}{|E_{\text{in}}|}.
\label{eq:innersupport}
\end{equation}
\item \textbf{判断低支持度并引入放宽肤色候选：}当 $\rho_{\text{in}} \ge 0.18$ 时，说明内椭圆中已有较多初始肤色像素，无需额外重选；当 $\rho_{\text{in}} < 0.18$ 时，判定该图内椭圆像素过少，计算阈值系数放宽至 $k_{\text{rec}}=2.00$ 时的肤色二值掩膜 $M_{\text{relax}}$，并截取其内椭圆部分 $M_{\text{relax}} \cap E_{\text{in}}$。
\item \textbf{闭运算连接与核心种子重叠过滤：}在 $E_{\text{in}}$ 内部，对 $M_{\text{relax}} \cap E_{\text{in}}$ 执行一次较大尺度的椭圆核闭运算（最终选用 $27\times27$ 核），把分散的脸部皮肤块连接成片，再与核心种子 $S$ 取并集：
\begin{equation}
A = \bigl(\operatorname{Close}_{27}(M_{\text{relax}}\cap E_{\text{in}})\cap E_{\text{in}}\bigr)\cup S.
\label{eq:allowmask}
\end{equation}
考虑到即使在内椭圆内部，放宽阈值后额头边缘也可能出现少量孤立的刘海误检块，我们对 $A$ 做 8 连通域分解，只保留与初始核心种子 $S$ 重叠像素数不少于 $24$ 的连通分量，最后在 $E_{\text{in}}$ 内部做一次漫水填孔，得到重建掩膜 $M_{\text{rec}}$：
\begin{equation}
M_{\text{rec}} = \operatorname{FillHoles}\!\left(\bigcup_{C_j\in\mathcal{C}_8(A):\,|C_j\cap S|\ge 24} C_j\right)\cap E_{\text{in}}.
\label{eq:lowsupport}
\end{equation}
\end{enumerate}

这样处理后，新增像素一方面被严格限制在内椭圆 $E_{\text{in}}$ 内部，另一方面必须通过放宽后的肤色闭运算与原有的高可信种子 $S$ 连成一体才会被保留，避免了不看图像肤色、直接把整个内椭圆涂满的做法。

\subsection{内外分区形态学操作与整体连接流程}

无论是常规样本还是触发了低支持重建的样本，最后都需要将内椭圆核心区与内外椭圆之间的过渡环带合成为完整的人脸掩膜。按第~\ref{sec:dual_ellipse}~章划定的三个空间区域，分区形态学连接的具体执行顺序如下：
\begin{enumerate}
\item \textbf{内椭圆核心区处理：}对保留的核心候选 $M_{\text{sel}}$ 执行 $11\times11$ 椭圆核闭运算并填孔，截取其落在内椭圆 $E_{\text{in}}$ 内的部分；若当前样本触发了低支持度重建（$\rho_{\text{in}}<0.18$），则将重建结果 $M_{\text{rec}}$ 一并并入核心区掩膜 $M_{\text{strong}}$：
\begin{equation}
M_{\text{strong}} = \bigl(\operatorname{FillHoles}(\operatorname{Close}_{11}(M_{\text{sel}}))\cap E_{\text{in}}\bigr) \cup M_{\text{rec}}.
\label{eq:mstrong}
\end{equation}
\item \textbf{过渡环带处理：}对 $M_{\text{sel}}$ 执行 $5\times5$ 椭圆核弱闭运算，仅截取其落在过渡环带 $E_{\text{out}}\setminus E_{\text{in}}$ 内的部分，记为 $M_{\text{weak}}$：
\begin{equation}
M_{\text{weak}} = \operatorname{Close}_{5}(M_{\text{sel}})\cap (E_{\text{out}}\setminus E_{\text{in}}).
\label{eq:mweak}
\end{equation}
\item \textbf{合并、种子连通过滤与最终填孔：}将核心区与过渡环带合并为候选掩膜 $M_{\text{prop}} = M_{\text{strong}} \cup M_{\text{weak}}$。为防止过渡环带中残留与面部主体断开的孤立碎块，再次对 $M_{\text{prop}}$ 提取 8 连通域，仅保留与核心种子 $S$ 相交的连通分量，最后做一次内部填孔并限制在 $E_{\text{out}}$ 内，得到最终分割掩膜 $M_{\text{final}}$：
\begin{equation}
M_{\text{final}} = \operatorname{FillHoles}\!\left(\bigcup_{C_j\in\mathcal{C}_8(M_{\text{prop}}):\,C_j\cap S\neq\emptyset} C_j\right)\cap E_{\text{out}}.
\label{eq:zonalmorph}
\end{equation}
\end{enumerate}

代码~\textbf{5-1} 列出了低支持度重建与内外分区形态学连接的核心实现。

\noindent\textbf{代码 5-1}\quad 低支持放宽阈值重建与双椭圆分区形态学连接（dual\_ellipse.py）
\begin{lstlisting}
def reconstruct_low_support(relaxed, seed, inner, kernel_size, minimum_seed_pixels):
    allowed = morphology(relaxed & inner, cv2.MORPH_CLOSE, kernel_size) & inner
    allowed |= seed & inner
    count, labels = cv2.connectedComponents(allowed, connectivity=8)
    retained = [c for c in range(1, count)
                if np.count_nonzero((labels == c) & seed.astype(bool)) >= minimum_seed_pixels]
    return fill_holes(np.isin(labels, retained).astype(np.uint8)) & inner

# 分区形态学连接：内椭圆强闭运算+低支持重建，过渡区弱闭运算，种子连通锚定
weak = morphology(selected, cv2.MORPH_CLOSE, config["weak_kernel"]) & outer
strong = fill_holes(morphology(selected, cv2.MORPH_CLOSE, config["strong_kernel"])) & inner
if relaxed is not None and inner_support < config["low_support_threshold"]:
    strong |= reconstruct_low_support(relaxed, seed, inner,
                                      config["reconstruction_close_kernel"],
                                      config["minimum_reconstruction_seed_pixels"])
proposal = strong | (weak & (outer & (1 - inner)))
count, labels = cv2.connectedComponents(proposal, connectivity=8)
anchor_labels = np.unique(labels[seed.astype(bool)])
final = fill_holes(np.isin(labels, anchor_labels[anchor_labels > 0]).astype(np.uint8)) & outer
\end{lstlisting}

\FloatBarrier
\section{颜色模型的能力边界与适用范围}\label{sec:boundary}

前五章分别处理了数据、统计、判决、独立测试与补全尝试。本章不引入新的实验，而是回到已经固定的结果，回答一个更基本的问题：这套以 Cr/Cb 阈值为核心的方法，究竟在什么意义上有效，又在什么意义上失效。本章用到的两项统计——F1 分布与精确率上限——均由已保存的逐图像计数和原始类别计数直接算出，没有重新预测，也没有调整参数；它们属于对固定结果的事后解释，不构成新一轮独立测试。

\subsection{方法的实际构成}\label{subsec:composition}

把最终流程拆开，只包含三个判断。第一，像素的 Cr、Cb 是否同时落在以训练均值为中心、以 $k$ 倍标准差为半宽的矩形内；第二，画面中面积最大的 8 连通前景块是否就是脸部；第三，不与图像边界连通的前景空洞是否属于脸部。三者分别对应颜色先验、单主体先验和闭合轮廓先验。补全只是把闭合先验做得更强，没有增加新的判别信息。

这三个先验都不涉及“脸长什么样”。方法没有使用形状、纹理、边缘或五官位置，并在颜色转换后主动丢弃了 Y 通道。因此，我们更愿意把它描述为“肤色检测加两个空间约束”，而不是“人脸分割”。后文对优劣的判断都建立在这一认识上：误差主要不来自阈值是否精细，而来自三个先验在具体图像上是否成立。

这种简化并非疏忽。丢弃 Y 是为了让判决只依赖色度，减少曝光变化的影响；采用最大连通域，是为了在不引入人脸检测器的前提下排除零散干扰；填孔则是为了在不使用五官标签的前提下恢复眼睛、眉毛和嘴部。三项选择都以“不引入额外模型”为代价，换取实现的简单与结果的可解释。这些代价在第~\ref{sec:test}~章与第~\ref{sec:completion}~章的失败样例中均有体现。

\subsection{三个工作点}

表~\ref{tab:points} 把三次判断的结果放在一起。纯阈值是颜色单独作用时的位置；初版加入区域与填孔；补全再加膨胀与腐蚀。

\begin{table}[H]
\centering\small
\setlength{\tabcolsep}{4pt}
\caption{三种方案在测试集上的位置与证据状态}\label{tab:points}
\begin{tabular}{lrrrrl}
\toprule
方案 & Precision & Recall & F1 & IoU & 证据状态 \\
\midrule
纯阈值 & 52.91 & 76.82 & 62.66 & 45.63 & 独立测试 \\
初版 & 66.43 & 75.55 & 70.69 & 54.67 & 独立测试 \\
补全后 & 63.63 & 79.55 & 70.70 & 54.68 & 开发后对比 \\
\bottomrule
\end{tabular}
\end{table}

初版把精确率从 52.91 拉到 66.43，召回几乎不动；补全把召回推到 79.55，精确率退回 63.63。初版与补全后落在同一条 F1 等值线附近；基线不在同一条线上，F1 低约 8 个百分点。图~\ref{fig:pr} 把这层关系画了出来。

\begin{figure}[H]
\centering
\begin{tikzpicture}
\begin{axis}[
    width=0.82\textwidth,
    height=0.60\textwidth,
    xlabel={召回率 $R$（\%）},
    ylabel={精确率 $P$（\%）},
    xmin=50, xmax=85,
    ymin=45, ymax=75,
    xtick={50,55,60,65,70,75,80,85},
    ytick={45,50,55,60,65,70,75},
    grid=major,
    grid style={draw=gray!20},
    legend style={at={(0.02,0.98)},anchor=north west,font=\small,draw=none,fill=none},
    tick label style={font=\small},
    label style={font=\small},
]
\addplot[domain=64:85,samples=80,dashed,draw=gray!70,thick,forget plot] {70.7*x/(2*x-70.7)};
\node[gray!80,font=\scriptsize,anchor=north west] at (axis cs:82,61.5) {$F1=0.707$};
\addplot[domain=70:85,samples=60,dashed,draw=gray!45,forget plot] {65*x/(2*x-65)};
\node[gray!60,font=\scriptsize,anchor=north west] at (axis cs:80,54.0) {$F1=0.65$};
\addplot[only marks,mark=square*,mark size=2.6pt,draw=blue!60!black,fill=blue!60!black] coordinates {(76.82,52.91)};
\addlegendentry{纯阈值}
\addplot[only marks,mark=*,mark size=2.8pt,draw=myteal,fill=myteal] coordinates {(75.55,66.43)};
\addlegendentry{初版（区域+填孔）}
\addplot[only marks,mark=o,mark size=3.2pt,draw=myteal,fill=white,mark options={line width=0.8pt}] coordinates {(79.55,63.63)};
\addlegendentry{补全后（开发后）}
\node[font=\scriptsize,anchor=south west] at (axis cs:75.55,66.43) {70.69};
\node[font=\scriptsize,anchor=north west] at (axis cs:79.55,63.63) {70.70};
\node[font=\scriptsize,anchor=north] at (axis cs:76.82,52.91) {62.66};
\end{axis}
\end{tikzpicture}
\caption{三种方案在测试集上的精确率—召回率位置。实心点为独立测试，空心点为开发后对比；虚线为 F1 等值线。}
\label{fig:pr}
\end{figure}

这张图说明，补全没有让模型“更准”，只是把判决边界向外推了一圈：被包进来的像素中只有 35.45\% 是脸。初版与补全后之间的选择，不是“哪个方法更好”，而是“愿意用多少精确率换多少召回”。

\subsection{有效性的来源}

从结果分布看，32 张测试图的初版 F1 并不呈连续分布。不低于 0.85 的有 11 张，其中 4 张超过 0.92；低于 0.2 的有 5 张；0.2 至 0.5 之间没有样本。补全后这个形状基本不变：11 张不低于 0.85，5 张低于 0.2，只有个别图像跨过区间边界。方法在多数图像上表现较好，在少数图像上几乎完全失败，中间状态很少。

在我们看来，这种两段式分布与第~\ref{subsec:composition}~节的结构假设一致。当脸部构成画面中占主导的肤色连通区域时，三个先验同时成立：颜色候选接近完整，最大连通域不会选错，填孔也能补回被包围的五官。当这一条件不成立时，错误会沿同一条链条被放大，而不是表现为分数的平缓下降。当前样本只有 32 张，这一形态也可能带有偶然性；但它至少提示，方法的成败更接近“条件是否满足”，而非“参数是否精细”。

在满足条件的图像上，方法的表现相当稳定。样本 00610 的原始候选为 80,328 像素，最大连通域保留 79,899 像素，填孔补回 9,490 个正确像素，最终 F1 为 0.930，召回率为 97.7\%。00585、00629、00538 等图的情况类似。这些图像的一个共同特征是正面或近正面、脸部占画面比例较大、光照较均匀，头发与背景的颜色和肤色可分。对这些图像，颜色候选本身已接近完整，后处理只需少量修补。

方法的第二个优势体现在对五官的补偿上。颜色规则本来不覆盖眼睛、眉毛和嘴部，但填孔把这一组的召回率从 32.6\% 提高到 68.3\%，补回 69,662 个正确像素；皮肤与鼻子组也从 66.4\% 提高到 76.3\%。这说明“脸是一个被肤色包围的连通区域”这一先验在实际图像中相当强。方法并不区分这些像素的语义类别，只是利用了它们被脸部轮廓包围这一拓扑事实。就本实验的目标而言，这一补偿是有效的；但它也解释了补偿的上限：只有被完整包围的区域才能被补回。

第三个优势是可解释与可复现。方法只有四个统计参数和一个整数 $k$，没有随机初始化，也没有在测试阶段重新估计任何量。每一次失败都可以定位到颜色候选、区域选择或填孔中的某一步：00609 的错误始于颜色候选，00718 的错误始于区域选择，00676 的误差则在填孔阶段被放大。这种可归因性使讨论不必停留在“效果不好”，而可以具体到“哪一条先验没有成立”。此外，方法不需要训练，单张图像只经历一次颜色变换和两次连通域操作，结构上比需要训练的方法轻。

\subsection{失效的来源}

方法的失效可以按来源分为四类：颜色、空间、语义，以及评价口径。前三类属于方法本身，第四类属于我们如何定义对错。表~\ref{tab:failure} 给出摘要。

\begin{table}[H]
\centering\small
\setlength{\tabcolsep}{4pt}
\caption{失效的四类来源}\label{tab:failure}
\begin{tabular}{p{1.5cm}p{4.0cm}p{6.0cm}p{3.0cm}}
\toprule
来源 & 不成立的假设 & 主要证据 & 调参能否解决 \\
\midrule
颜色 & 肤色落在训练矩形内 & 00609：候选 5,685 像素、56 块、召回 1.6\%；低亮度组 F1 下降 3.8 个百分点 & 部分缓解，误检会回升 \\
空间 & 存在占主导的脸部连通块 & 00718：最大块 58,716 像素，TP 仅 2,705；脖子占误检 23.4\% & 需更换区域评分 \\
语义 & 肤色像素都属于脸 & 脖子与双耳 37.4 万像素；精确率上限约 85\% & 不能 \\
评价 & 整脸真值与候选边界一致 & 逐图 Mask IoU 0.572，Boundary IoU 0.233 & 不能，属定义问题 \\
\bottomrule
\end{tabular}
\end{table}

颜色层。当图像存在明显偏色或照度不足时，真实肤色的色度会偏离训练分布，$k=1.25$ 的矩形随之把大量正确像素排除在外。00609 是典型：$k=1.25$ 下只剩 5,685 个候选像素，分成 56 个连通块，最大块 1,063 像素，最终召回率 1.6\%。亮度分组的结果与此一致：较低亮度组的初版 F1 为 52.8\%，比基线低 3.8 个百分点，正确像素减少约 9.6 万，漏检增加约 9.6 万。补全对这类图几乎没有作用：00609 的召回只从 1.59\% 到 1.78\%。需要说明，亮度分组是事后按真值区域平均 $Y$ 划分的，只有 11 张图，且亮度与场景、肤色、妆容混杂，因此这里只能说明低亮度条件下风险更高，不能把下降单独归因于亮度。

空间层。最大连通域假设画面中存在一个占主导地位的脸部候选。当与肤色相近的头发形成更大连通块时，这一假设直接失效。00718 中，棕色头发构成 58,716 像素的最大块，其中正确像素只有 2,705 个，误检 56,011 个；填孔随后又增加 110 个正确像素和 10,871 个误检像素，最终 F1 为 0.045。补全把召回从 5.12\% 推到 9.78\%，仍然无济于事；00578 中，补全甚至沿着头发继续扩张，新增 22,692 个误检像素。这个例子说明，面积只能表示区域大小，不能表示区域属于什么。类似地，脖子与脸部通常连成同一块，无法通过连通性切开；在最终误检中，头发与脖子合计占 84.0\%，其中脖子单独占 23.4\%。

语义层。这是方法的真正上限。耳朵和脖子本身也是皮肤，颜色规则没有理由排除它们。测试集中，脖子与双耳合计约 37.4 万像素，占全部有效像素的 4.5\%，而脸部正类约 215 万像素。即使颜色判决完全正确、不再产生任何头发或背景误检，只要结构仍是“最大连通域加填孔”，这部分像素仍会被计入前景，精确率的上限约为 85\%。当前初版的精确率为 66.4\%，与这一上限之间还有约 19 个百分点，说明头发与背景误检仍是主要矛盾；但即使把这些误检全部消除，纯颜色模型也无法达到更高的精确率。这一计算建立在“所有肤色像素都被判为正类”的假设上，实际方法可能因连通性排除部分耳朵，因此 85\% 是量级上的上限，而非严格界。它的意义在于指出：剩余误差中有一部分不是参数问题，而是“颜色”这一信息本身不包含“脸”这一语义。

评价口径。还有一层不确定性来自评价本身。整脸真值与颜色候选之间的错位是任务定义的一部分：眼睛、眉毛和嘴部被计入正类，却不参与颜色建模。边界质量也受到这一错位的影响。初版方案的逐图 Mask IoU 平均为 0.572，而 Boundary IoU 平均仅为 0.233；前者说明区域大体重合，后者说明轮廓并不可信。填孔产生的是由外围轮廓决定的平滑边界，而不是真实的发际线或下颌线；标注本身在这些位置也存在歧义。因此，本方法适合给出粗略区域，不适合用于需要精确边缘的场合。

\subsection{适用条件}

把上述结果整理成条件，可以得到一个相对清晰的适用范围。这些条件来自 32 张测试图的观察，属于经验总结，不是经过独立验证的判据。表~\ref{tab:applicability} 给出摘要。

\begin{table}[H]
\centering\small
\setlength{\tabcolsep}{4pt}
\caption{适用条件与实验依据}\label{tab:applicability}
\begin{tabular}{p{4.2cm}p{7.2cm}p{3.4cm}}
\toprule
场景条件 & 实验观察 & 判断 \\
\midrule
单人、近正面、脸部占比较大、光照均匀 & 11 张 F1 不低于 0.85，4 张超过 0.92 & 适用 \\
偏色或照度不足 & 00609 F1 为 0.031；低亮度组下降 3.8 个百分点 & 风险高，需预处理 \\
头发或背景与肤色接近 & 头发占误检 60.6\%；00718 F1 为 0.045 & 易失效 \\
脸与脖子相连 & 脖子占误检 23.4\% & 精度受限 \\
多人、侧脸、遮挡 & 本数据未覆盖 & 未验证，不宜外推 \\
需要精确边缘 & Boundary IoU 0.233 & 不适用 \\
\bottomrule
\end{tabular}
\end{table}

据此，我们把适用场景限定为：单人、正面或近正面、脸部在画面中占较大比例、光照较均匀且与背景有较明显色差的图像；任务应是粗粒度区域提取或前置筛选，而不是精确边缘。多人、侧脸、遮挡等情形在本数据中没有覆盖，我们不把它们列入适用范围，也不据此判断方法在那些情形下的表现。概括而言，这套方法是一个在受控条件下可用的肤色先验，其适用范围主要由拍摄条件决定，参数调节只能在其内部微调。

如果必须在初版与补全后之间选一个，判断依据不是 F1，而是任务更在意什么。以 F1 和边界为准，选初版；以覆盖完整为准，选补全后。两者的 F1 相同，差别在精确率与召回的分配；这也是第~\ref{sec:completion}~章反复出现的“工作点”一词的实际含义。

\subsection{与其他颜色方法的关系}

最后，把本方法放回颜色分割的谱系中，可以更清楚地看出它的位置。

基于 RGB 的规则更简单，但三个通道都混有亮度信息，阈值同时受光照强度与颜色影响；Cr/Cb 把亮度分离出去，代价是丢弃了 Y 中的明暗线索。本实验没有直接比较两者，00609 的结果只能说明色度通道同样会受白平衡影响，不能说明 RGB 规则一定更差。HSV 的色调对光照变化相对稳健，但皮肤通常落在低饱和区域，色调在低饱和时并不稳定。

训练分布本身还提供了另一条线索。Cr 与 Cb 的相关系数约为 $-0.76$，联合分布呈倾斜的椭圆，而当前判决使用的是轴对齐矩形。矩形的四个角会包含肤色分布之外的取值组合，而倾斜方向上的肤色又可能被矩形的边切掉。若改用马氏距离或椭圆边界，不必更换颜色空间，只需在现有均值、标准差之外增加一个相关系数，把“各通道独立”的假设换成“两通道联合”的假设。就我们看到的分布形态而言，这是改动最小、针对性最直接的一处。本实验没有实现这一方案，因此不能给出它可能带来的提升；这里只把它作为由数据直接支持的一个判断记录下来。

逐图自适应阈值是另一种常见思路。它可以根据每张图的色度分布重新确定范围，对偏色更稳健；但在本任务中，背景中的木质墙面、暖色衣物同样会形成大面积色度聚集，自适应阈值有可能把这些区域纳入候选，与最大连通域的假设冲突。是否采用，取决于场景中是否存在稳定的肤色参照，而不能只按“能否适应光照”来判断。

至于基于学习的分割方法，它们能够利用形状、纹理和语义信息，理论上可以越过语义层的限制。但这类方法需要标注数据与训练过程，其结果也不再能简单归因到“颜色范围是否合适”。对于本实验提出的问题——颜色本身能走多远——引入学习模型后，这一问题就不再被回答。因此，我们选择在颜色框架内把边界刻画清楚，而不是用更强的模型换取更高的分数。

\subsection{小结}

本章的结论可以概括为一句话：这套方法的精度上限由它看不见的信息决定，而不是由 $k$ 或核大小决定。颜色先验、单主体先验和闭合轮廓先验在受控条件下同时成立，方法就能达到 0.9 左右；其中任何一条不成立，误差就会沿流程被放大，而不是被后处理吸收。参数调整只能改变三条先验的松紧，不能补上它们缺少的形状与语义。就本实验而言，它的价值主要在于实现代价低、结果可解释、失败可归因，而不在精度上限。

\FloatBarrier
\section{实验总结与展望}\label{sec:conclusion}

本报告基于 162 张带像素级标签的人脸图像，完整呈现了从实验三（YCrCb 颜色建模与基础连通域后处理）到实验四（双眼定位引导的双椭圆空间约束与低支持度放宽阈值重建）的递进研究过程。全文的核心工作与结论可总结为以下四个方面：

第一，\textbf{YCrCb 色度建模与矩形阈值的覆盖特性：}训练集 638.8 万个皮肤与鼻子像素在 Cr/Cb 平面上呈明显的倾斜椭圆聚集（两通道标准差均约为 9.4，相关系数约为 $-0.7644$）。实验三选定的 $k=1.25$ 矩形阈值侧重于抑制非肤色背景，实质上截取了高密度肤色核心区；这一特性保证了常规光照下的高精确率，但在暗光或冷偏色图像上容易导致面部候选碎裂。

第二，\textbf{基础连通域后处理的阶段增益：}在实验三中，$k=1.25$ 阈值配合最大 8 连通域筛选与内部空洞填充，在 32 张测试图上取得全局 F1 70.69\%、IoU 54.67\%，比纯阈值基线分别提升 8.03 和 9.04 个百分点。分阶段像素分解表明，阈值收紧、最大连通域筛选与内部填孔的正负像素交换率分别为 $1:1.6$、$1:8.3$ 和 $3.3:1$，三者协同完成了背景去噪与封闭五官恢复。

第三，\textbf{纯颜色与无边界形态学操作的等值线瓶颈：}误差语义分解显示，头发（60.58\%）与脖颈（23.40\%）合计占实验三最终误检的 83.98\%。在缺乏人脸几何边界的情况下，$41\times41$ 椭圆核“膨胀—填孔—腐蚀”补全虽将测试集召回率提高 4.00 个百分点（达 79.55\%），但新增前景像素中 64.55\% 属于外围头发与背景误检，精确率下降 2.80 个百分点，全局 F1（70.70\%）保持不变。这证明无空间锚点的形态学操作只能在 F1 等值线上调节精确率与召回率的分配。

第四，\textbf{双眼引导双椭圆约束与低支持重建的协同突破：}实验四通过分阶段连通域诊断与轮廓直接拟合实验，证实了待修整肤色轮廓本身易受粘连头发或局部碎片带偏；进而利用灰度黑帽变换结合训练集左右眼独立 $3\sigma$ 窗口先验定位双眼（验证集 $e\le 0.25$ 准确率达 28/32，配合严格/放宽两级面积比回退覆盖 31/32 张图，无配对样本自动回退 Lab3），在双眼局部坐标系下生成正中及左右侧向候选（$\pm 0.16\,d_{\text{eye}}$，偏移门槛 $\Delta S_{\text{lat}}=0.06$）与内外嵌套双椭圆（$\lambda_{\text{in}}=0.82,\;\lambda_{\text{out}}=1.20$）。外椭圆有效截断了大面积粘连的头发与脖颈误检，内椭圆针对支持率 $\rho_{\text{in}}<0.18$ 的碎裂样本触发 $k_{\text{rec}}=2.00$ 放宽肤色重选与 $27\times27$ 闭运算受限重建，最后经内椭圆 $11\times11$ 强闭运算与过渡带 $5\times5$ 弱闭运算完成核心种子锚定连接。在 32 张验证图上，实验四最终方案将全局 Precision 从 64.86\% 提升至 79.88\%（+15.02 个百分点），Recall 从 82.70\% 提升至 86.15\%（+3.45 个百分点），全局 F1 从 72.70\% 跃升至 82.89\%（+10.19 个百分点），IoU 从 57.11\% 提升至 70.79\%（+13.68 个百分点），实现了精确率与召回率的同步跃升。

展望后续工作，可在现有颜色—几何混合框架下沿三个方向进一步拓展：一是将 Cr/Cb 轴对齐矩形阈值升级为考虑协方差结构的马氏距离椭圆判别；二是结合鼻部与嘴部几何线索优化侧向椭圆姿态选择准则，充分释放侧脸修正潜力（理想姿态上界 F1 达 84.44\%）；三是在内椭圆上缘引入局部边缘梯度与纹理约束，进一步抑制暖调刘海在放宽阈值重建时的误检混入。

\FloatBarrier

\begin{thebibliography}{9}
\bibitem{hsu2002face}
Hsu R L, Abdel-Mottaleb M, Jain A K.
Face detection in color images[J].
IEEE Transactions on Pattern Analysis and Machine Intelligence, 2002, 24(5): 696--706.

\bibitem{vezhnevets2003survey}
Vezhnevets V, Sazonov V, Andreeva A.
A survey on pixel-based skin color detection techniques[C].
Proceedings of GraphiCon, 2003: 85--92.

\bibitem{kovac2003skin}
Kovac J, Peer P, Solina F.
Human skin color clustering for face detection[C].
Proceedings of IEEE EUROCON, 2003, 2: 144--148.

\bibitem{cheng2021boundary}
Cheng B, Girshick R, Doll\'{a}r P, et al.
Boundary IoU: Improving object-centric image segmentation evaluation[C].
Proceedings of the IEEE/CVF Conference on Computer Vision and Pattern Recognition, 2021: 15334--15342.
\end{thebibliography}
\end{document}
